\subsection{The case of abelian groups}

Let G be a locally compact group, \(N \subset \Gamma^{0}\) the subspace of normalized Haar measures on the discrete subgroups of G, and \(N_{c}\) the subset of N corresponding to the discrete subgroups H of G such that G/H is compact; thus \(N_{c} = N \cap \Gamma_{c}^{0}\) in the notations of No.~3, Prop. 6; and if the group G is generated by a compact neighborhood of e, it follows from No.~3, Prop. 6 that \(N_{c}\) is open in N (hence is locally compact by No.~3, Cor. of Prop. 7) and that the restriction to \(N_{c}\) of the mapping \(\alpha \mapsto \|\mu_{\alpha}\|\) is vaguely continuous.

PROPOSITION 9. --- Let G be a locally compact abelian group, generated by a compact neighborhood of e. For a subset A of \(N_{c}\) to be relatively compact in \(N_{c}\) , it is necessary and sufficient that it satisfy the following two conditions:

\begin{enumerate}
\def\labelenumi{(\roman{enumi})}
\item
  There exists an open neighborhood \(\mathbf{U}\) of \(e\) in \(\mathbf{G}\) such that \(\mathrm{H}_{\alpha} \cap \mathrm{U} = \{e\}\) for all \(\alpha \in \mathbf{A}\) .
\item
  There exists a constant \(k\) such that \(\mu_{\alpha}(\mathrm{G} / \mathrm{H}_{\alpha}) \leqslant k\) for all \(\alpha \in \mathbf{A}\) . If \(\mathbf{A} \subset \mathbf{N}_c\) is relatively compact in \(\mathbf{N}_c\) , it is a fortiori so in \(\mathbf{N}\) , and the necessity of the conditions (i) and (ii) therefore follows from No.~3, Props. 6 and 7 (without assuming \(\mathbf{G}\) to be abelian). Conversely, suppose that \(\mathbf{A} \subset \mathbf{N}_c\) satisfies these conditions; if \(\overline{\mathbf{A}}\) is the closure of \(\mathbf{A}\) in \(\mathbf{N}\) , then \(\overline{\mathbf{A}}\) is compact by virtue of No.~3, Prop. 7; moreover, since \(\alpha \mapsto \| \mu_{\alpha} \|\) is lower semi-continuous on \(\Gamma^0\) for the vague topology (No.~2, Prop. 4), the condition (ii) implies that one also has \(\| \mu_{\alpha} \| \leqslant k\) for all \(\alpha \in \overline{\mathbf{A}}\) . Now, since \(\mathbf{G}\) is abelian, \(\mu_{\alpha} = \mu / \alpha\) is a Haar measure on the group \(G / H_{\alpha}\) , and \(G / H_{\alpha}\) is therefore compact for every \(\alpha \in \overline{\mathbf{A}}\) (Ch. VII, §1, No.~2, Prop. 2). This means that \(\overline{\mathbf{A}} \subset N_c\) , thus \(\mathbf{A}\) is relatively compact in \(N_c\) .
\end{enumerate}

COROLLARY. --- Let G be a locally compact abelian group, generated by a compact neighborhood of e and satisfying the condition (L) of No.~4.

Let \(D_c\) be the set of discrete subgroups H of G such that G/H is compact, and, for every \(H \in D_c\) , let \(v(H)\) be the total mass \(\mu_\alpha(G/H)\) , where \(\mu_\alpha\) is the quotient measure of \(\mu\) by the normalized Haar measure \(\alpha\) of H. For a subset A of the space \(D_c\) to be relatively compact in \(D_c\) , it is necessary and sufficient that it satisfy the following two conditions:

\begin{enumerate}
\def\labelenumi{(\roman{enumi})}
\item
  There exists an open neighborhood U of e in G such that \(\mathbf{H} \cap \mathbf{U} = \{e\}\) for all \(\mathbf{H} \in \mathbf{A}\) .
\item
  There exists a constant \(k\) such that \(v(\mathbf{H}) \leqslant k\) for all \(\mathbf{H} \in \mathbf{A}\) .
\end{enumerate}

Taking into account Prop. 9, this follows at once from the fact that \(D_{c}\) is the image of \(N_{c}\) under the canonical bijection of N onto D, and the fact that, under the hypotheses made, this bijection is a homeomorphism (No.~4, Prop. 8).

Example. --- Let us take \(G = R^{n}\) and for \(\mu\) the Lebesgue measure; all of the hypotheses of the Cor. of Prop. 9 are satisfied. The discrete subgroups H of G such that G/H is compact are none other than the discrete subgroups of rank n (GT, VII, §1, No.~1, Th. 1); such a subgroup H is generated by a basis \((a_{i})_{1 \leqslant i \leqslant n}\) of \(R^{n}\) , and

\[
v (\mathbf {H}) = | \det (a _ {1}, \dots , a _ {n}) |
\]

(the determinant being taken with respect to the canonical basis of \(\mathbf{R}^n\) ) (Ch. VII, §2, No.~10, Th. 4). The space \(\mathbf{D}_c\) can be interpreted here in the following way: every subgroup \(\mathbf{H} \in \mathbf{D}_c\) is the transform \(g \cdot \mathbf{Z}^n\) of the subgroup \(\mathbf{Z}^n\) by an element \(g \in \mathbf{GL}(n, \mathbf{R})\) , and the subgroup of \(\mathbf{GL}(n, \mathbf{R})\) leaving \(\mathbf{Z}^n\) stable may be identified with \(\mathbf{GL}(n, \mathbf{Z})\) . Consequently \(\mathbf{D}_c\) may be canonically identified, as a (non-topological) homogeneous space, with \(\mathbf{GL}(n, \mathbf{R}) / \mathbf{GL}(n, \mathbf{Z})\) . On the other hand, \(\mathbf{GL}(n, \mathbf{R})\) operates continuously in \(\mathbf{R}^n\) , hence also in \(\mathcal{M}_{+}(\mathbf{R}^{n})\) for the vague topology (§3, No.~3, Prop. 13), hence in the subspace \(\mathrm{N}_c\) of \(\mathcal{M}_{+}(\mathbf{R}^{n})\) ; moreover, the canonical homeomorphism (No.~4, Prop. 8) of \(\mathrm{N}_c\) onto \(\mathrm{D}_c\) is compatible with the laws of operation of \(\mathbf{GL}(n, \mathbf{R})\) . Since \(\mathbf{GL}(n, \mathbf{R})\) is countable at infinity and \(\mathrm{D}_c\) is locally compact, the bijection of \(\mathbf{GL}(n, \mathbf{R}) / \mathbf{GL}(n, \mathbf{Z})\) onto \(\mathrm{D}_c\) defined above is a homeomorphism (Ch. VII, App. I, Lemma 2). The Cor. of Prop. 9 therefore gives a criterion for compactness in the homogeneous space \(\mathbf{GL}(n, \mathbf{R}) / \mathbf{GL}(n, \mathbf{Z})\) .

